The delivery is one folder per model, eight in all, each holding three frame directories of
thirteen Parquet files:

\begin{verbatim}
  <model>/
    I_mean/                 predictions_calibration_<ts>_<NN>.parquet  (13 files)
    II_q95/                 "
    III_conditional_mean/   "
\end{verbatim}

\noindent That is $8 \times 3 \times 13 = 312$ files and a little over 180~MB. The timestamp in each
filename is the model's own run, so it is the same within a model and differs between models; the
two-digit suffix \texttt{\_00} to \texttt{\_12} is the forecast origin.

Parquet is a columnar table format. It is not a text format and will not open in a text editor,
but every analysis environment reads it directly, with no conversion step:

\begin{quote}
\small
\textbf{Python} \\
\texttt{import pandas as pd} \\
\texttt{f = "purple\_alien/I\_mean/predictions\_calibration\_20260928\_201214\_00.parquet"} \\
\texttt{df = pd.read\_parquet(f)}

\medskip
\textbf{R} \\
\texttt{library(arrow)} \\
\texttt{f <- "purple\_alien/I\_mean/predictions\_calibration\_20260928\_201214\_00.parquet"} \\
\texttt{df <- read\_parquet(f)}
\end{quote}

\subsection*{Thirteen Origins, and Why They Must Not Be Stacked}

The calibration test window is months 457 to 504 --- January 2018 to December 2021. Each model
forecasts 36 months ahead from each of thirteen successive starting points, one month apart.

\begin{table}[h!]
\centering
\small
\renewcommand{\arraystretch}{1.2}
\caption{The thirteen forecast origins. Each file covers 36 months from its own origin.}
\label{tab:origins}
\begin{tabular}{@{}llll@{}}
\toprule
\textbf{File} & \textbf{Forecast from} & \textbf{Months covered} & \textbf{Calendar} \\ \midrule
\texttt{\_00} & month 456 & 457--492 & Jan 2018 -- Dec 2020 \\
\texttt{\_01} & month 457 & 458--493 & Feb 2018 -- Jan 2021 \\
\texttt{\_02} & month 458 & 459--494 & Mar 2018 -- Feb 2021 \\
\multicolumn{4}{@{}l@{}}{\quad\textit{\ldots one month later at each step \ldots}} \\
\texttt{\_11} & month 467 & 468--503 & Dec 2018 -- Nov 2021 \\
\texttt{\_12} & month 468 & 469--504 & Jan 2019 -- Dec 2021 \\
\bottomrule
\end{tabular}
\end{table}

\noindent \textbf{The windows overlap, so the same cell-month appears in many files.} Months 469
to 492 --- twenty-four of the forty-eight in the span --- appear in \emph{all thirteen} origins,
each time at a different forecast age. Concatenating the thirteen
files and treating the result as a panel produces a table with up to thirteen different values for
the same cell and month, which is not a forecast of anything. Pick one origin, or carry the origin
as an explicit dimension. The audit notebook in the folder prints the overlap as a table if it
helps to see it.

\subsection*{Schema}

Every one of the 312 files holds exactly five columns and \textbf{2,333,448 rows} --- 64,818 cells
$\times$ 36 months, a complete block with no gaps. The schema is byte-identical across all 312.

\begin{table}[h!]
\centering
\small
\renewcommand{\arraystretch}{1.2}
\caption{The five columns, identical in every file.}
\label{tab:columns}
\begin{tabular}{@{}p{4.2cm}p{2.0cm}p{8.0cm}@{}}
\toprule
\textbf{Column} & \textbf{Type} & \textbf{Meaning} \\ \midrule
\texttt{month\_id}           & integer & VIEWS month index; 457 is January 2018 \\
\texttt{priogrid\_id}        & integer & PRIO-GRID cell identifier (see below) \\
\texttt{pred\_lr\_sb\_best}  & double  & predicted state-based fatalities \\
\texttt{pred\_lr\_ns\_best}  & double  & predicted non-state fatalities \\
\texttt{pred\_lr\_os\_best}  & double  & predicted one-sided fatalities \\
\bottomrule
\end{tabular}
\end{table}

\noindent \textbf{\texttt{priogrid\_id} is a cell, not a coordinate.} PRIO-GRID is a $720 \times
360$ grid of half-degree cells covering the globe, numbered from 1 in the south-west corner and
running east before north. To place a row on a map:

\begin{table}[h!]
\centering
\small
\renewcommand{\arraystretch}{1.3}
\caption{Converting \texttt{priogrid\_id} to coordinates. The result is the cell centre.}
\label{tab:priogrid}
\begin{tabular}{@{}p{4.6cm}p{9.4cm}@{}}
\toprule
Row (south to north) & $\texttt{row} = \lfloor (\texttt{priogrid\_id} - 1) / 720 \rfloor$ \\
Column (west to east) & $\texttt{col} = (\texttt{priogrid\_id} - 1) \bmod 720$ \\
Latitude  & $-90 + (\texttt{row} + 0.5) \times 0.5$ \\
Longitude & $-180 + (\texttt{col} + 0.5) \times 0.5$ \\
\bottomrule
\end{tabular}
\end{table}

\noindent In R:

\begin{quote}
\small
\texttt{d\$lat <- -90  + (((d\$priogrid\_id - 1) \%/\% 720) + 0.5) * 0.5} \\
\texttt{d\$lon <- -180 + (((d\$priogrid\_id - 1) \%\%  720) + 0.5) * 0.5}
\end{quote}

\noindent As a check: in \texttt{purple\_alien}'s first origin the largest predicted state-based
value is 58.0 fatalities in month 462 --- June 2018 --- at \texttt{priogrid\_id} 181875, which
resolves to $36.25^\circ$N, $37.25^\circ$E, in northern Syria. The 64,818 cells delivered here are
the land subset of the 259,200 in the full grid; their centres span $55.75^\circ$S to
$83.25^\circ$N.

Values are fatality counts. They are continuous rather than integer, because a mean over draws is
not a whole number, and they are on the original count scale --- not logged, not scaled, not
truncated. Observed values are not included (Section~\ref{sec:excluded}).

\subsection*{Checking the Delivery Without Running Anything}

The folder contains \texttt{audit\_selector\_frames.ipynb}, which audits the folder it sits in. It
opens all 312 files, checks the columns and types, checks for non-finite, negative and duplicated
values, verifies that every file is a complete $64{,}818 \times 36$ block, draws the global maps
and the per-step profiles, and ends in a cell that \emph{raises} rather than prints if anything is
wrong.

\textbf{Its outputs are saved from the run that checked this exact copy of the data}, so it can be
read without executing it. To re-run it, start Jupyter in the folder and Run All; it needs only
\texttt{pandas}, \texttt{pyarrow}, \texttt{numpy} and \texttt{matplotlib}, and reads the model
directories beside it. Nothing else in this delivery needs to be verified by hand, and there is no
checksum file to run.

For the record, rather than as something to take on trust: across all 312 files and all 2.18 billion
delivered values there is \textbf{no missing, negative or infinite value}, every file is a complete
$64{,}818 \times 36$ block with no repeated \texttt{(month\_id, priogrid\_id)} pair, and the
columns and types are exact in every file. Two further checks were run for this note. The three
frames are written in the same row order, so they can be compared row by row without a join. And on
every one of the 242,678,592 rows, in all three targets, frame II and frame III are both greater
than or equal to frame I --- no exceptions, worst violation exactly zero. The second of those is
the one worth having: since $\text{mean} = P(y>0) \times \mathrm{E}[y \mid y>0]$, a row where the
conditional mean fell below the mean would mean the two frames had been collapsed from different
draws. None did.

One gap to record: the delivered folder carries \textbf{no per-model manifest.} The engine version,
run host, runtime and configuration commit for each model exist, in the run records held in
views-models, and the commit is on the title page of this note --- but they do not travel inside
the folder. A later delivery of this kind should ship them beside the data.
